\documentclass[10pt,a4paper]{amsart}
\usepackage[margin=19mm]{geometry}
\usepackage{amsmath,amssymb,mathtools}
\usepackage[scr=rsfs,cal=euler]{mathalfa}
\usepackage{xfrac}
\usepackage[unicode,hidelinks,bookmarks=false]{hyperref}
\newcommand{\ie}{i.e.\ }
\newcommand{\cf}{cf.\ }
\newcommand{\resp}{respectively\ }
\newcommand{\Subsection}{\subsection}
\providecommand{\nobreakdash}{\nobreak\hbox{-}}
\setlength{\emergencystretch}{2em}
\multlinegap0pt
\usepackage{bm}
\newcommand{\reva}[1]{#1}%{\textcolor{teal}{#1}}
\theoremstyle{plain}
\newtheorem{theorem}{Theorem}
\newtheorem{proposition}[theorem]{Proposition}
\newtheorem{lemma}[theorem]{Lemma}
\newtheorem{corollary}[theorem]{Corollary}
\theoremstyle{definition}
\newtheorem{definition}[theorem]{Definition}
\newtheorem{assumption}[theorem]{Assumption}
\theoremstyle{remark}
\newtheorem{remark}[theorem]{Remark}
\title[Order jumps in Spectral Deferred Corrections]{Order jumps in Spectral Deferred Corrections: Theorem 2}
\author{Eugen Bronasco, Joscha Fregin, Daniel Ruprecht, Gilles Vilmart}
\date{Source: arXiv:2604.03013v2 (CC BY 4.0)}
\begin{document}
\maketitle

\noindent\emph{Source and licence.} The delivered work is the article shown in the title above, version arXiv:2604.03013v2 (CC BY 4.0), distributed under the Creative Commons Attribution 4.0 International licence, as recorded in the arXiv licence metadata for this exact version and shown on the abstract page saved with this item. The full licence notice, which carries the licence URL and the address of that abstract page, is the bundled file \texttt{licenses/arxiv-2604.03013-CC-BY-4.0.txt}.\par\smallskip

\noindent\textit{Editorial scope.} Counted unit: the article's Theorem 2 (source0.tex lines 2052--2062, source label \texttt{thm:order\_jump}); statement and proof are the article's own text line for line, in the article's own notation ($\bm{A}^k_\Delta$, $C_{\eta_k}(W_k)$, $D_{Y_k}(\zeta_k)$, $p_K$, $\alpha^{[K]}(\tau)$, $\omega_{K+1}$), with no line omitted, substituted or re-flowed. The article's own revision-marking macro \texttt{\textbackslash reva}, which the retained lines use and which the article defines as the identity, is restated unchanged. Required context retained uncounted: the article's per-sweep order theorem (source0.tex lines 1905--1913), and its general-form and eta-form lemmas, which the proof invokes (lines 2280--2293 and 2253--2261); the proofs of those lemmas lie outside every retained span and are not delivered. Editorial formatting only: an AMS \texttt{amsart} preamble that restates that macro and the theorem environments the retained lines use, with \texttt{bm} loaded for the article's bold math; the article's Springer Nature class file \texttt{sn-jnl.cls} and its \texttt{planarforest.sty} are neither shipped nor input; sequential numbering of the retained environments in order of appearance, whereas the article prints them as Theorem 2 and the supporting lemmas of its sections 4--6. No citation occurs inside any retained line, so the wrapper carries no bibliography. No asset-inclusion command, no drawing environment and no image asset occurs in this package.


\section*{Retained context: thm:order\_per\_sweep (source0.tex lines 1905--1913)}

\begin{theorem}\label{thm:order_per_sweep}
    Let the internal stages of an SDC method $(\bm{A}_{\Delta}^0, \ldots, \bm{A}_{\Delta}^K, \bm{A}, \bm{b})$ be of height order $h_K - 1$.
    It then holds that
    
    \begin{enumerate}
        \item the method has height order $h_K$,
        \item the internal stages of $(\bm{A}_{\Delta}^0, \ldots, \bm{A}_{\Delta}^K, \bm{A}_{\Delta}^{K+1}, \bm{A}, \bm{b}, \bm{c})$ are of height order $h_K$ for any $\bm{A}^{K+1}_\Delta$.
    \end{enumerate}
\end{theorem}


\section*{Retained context: lemma:ai\_eta\_form (source0.tex lines 2253--2261)}

\begin{lemma}\label{lemma:ai_eta_form}
    Consider an SDC with $K$ iterations \reva{approximating a Runge--Kutta method satisfying $C(\eta)$. Given $\eta_k \in \mathbb{R}$ and $W_k \in \mathbb{R}^{\eta_k}$, assume} $\bm{A}_\Delta^k$ satisfies $C_{W_k}(\eta_k)$ with
    \[ 1 \leq \eta_k \leq \eta, \quad \eta_{k+1} \leq \eta_k + 1, \quad \text{for } k = 1, \dots, K.\]
    Then, given a tree $\tau$ with $|\tau| \leq \eta_K$, we have,
    \begin{equation}\label{eq:ai_eta_form}
        \alpha^{[K,i]} (\tau) = c_i^{|\tau|} \omega_K (\tau),
    \end{equation}
    \reva{for some constant $\omega_K (\tau)$ which depends on the tree $\tau$ and $W_k$ for $k=1,\dots,K$.}
\end{lemma}


\section*{Retained context: lemma:ai\_gen\_form (source0.tex lines 2280--2293)}

\begin{lemma}\label{lemma:ai_gen_form}
    Consider SDC with $K$ iterations which approximates a \reva{Runge--Kutta method of order $p$ satisfying $B(p), C(\eta), D(\zeta)$ with $p = \eta + \zeta$. Assume there exists a sequence of natural numbers $\eta_k \in \mathbb{N}$ and constants $W_k \in \mathbb{R}^{\eta_k}$ for $k=1,\dots, K$ such that $\bm{A}^k_\Delta$ satisfies $C_{\eta_k}(W_k)$ with
    \[ 1 \leq \eta_k \leq \eta \,, \quad \eta_{k+1} = \eta_k + 1 \,. \]
    Then, given $\tau = \{\tau_1 \cdots \tau_n\}$, we have,
    \begin{equation}\label{eq:ai_gen_form}
        \alpha^{[K]} (\tau) = \frac{1}{|\tau|} \omega_K (\tau_1) \cdots \omega_K (\tau_n),
    \end{equation}
    if one of the following conditions is satisfied:
    \begin{enumerate}
        \item $|\tau| \leq \eta_K + 1$,
        \item $\eta_K + 1 < |\tau| \leq \eta_K + \zeta_K + 1$ and $\bm{A}_\Delta^k$ satisfies $D_{Y_k}(\zeta_k)$ with $Y_k \in \mathbb{R}$ where,
            \[ 1 \leq \zeta_k \leq \zeta, \quad \zeta_k \leq \eta_k + 1 \,, \quad \text{and} \quad k = 1, \dots, K \,. \]
    \end{enumerate}}
\end{lemma}


\section{The counted unit: Theorem 2 and its complete original proof}

\begin{theorem}
    \label{thm:order_jump}
    Consider SDC with $K+1$ iterations which approximates a \reva{Runge--Kutta method of order $p$ satisfying $B(p), C(\eta), D(\zeta)$ with $p = \eta + \zeta$. Let the $K^{th}$ iteration of SDC be of iteration order $p_K$. Assume there exists a sequence of natural numbers $\eta_k \in \mathbb{N}$ and constants $W_k \in \mathbb{R}^{\eta_k}$ for $k=1,\dots, K+1$ such that $\bm{A}^k_\Delta$ satisfies $C_{\eta_k}(W_k)$ with
    \[ W_{K+1,p_K + 1} = \frac{1}{p_K + 1} \,, \quad \text{and} \quad 1 \leq \eta_k \leq \eta \,, \quad \eta_{k+1} = \eta_k + 1 \,. \]
    We have $p_{K+1} = p_K + 2$ if one of the following conditions is satisfied:
    \begin{enumerate}
        \item $p_K < \eta_{K+1}$,
        \item $\eta_{K+1} \leq p_K < \eta_{K+1} + \zeta_{K+1}$ and $\bm{A}_\Delta^k$ satisfies $D_{Y_k}(\zeta_k)$ for $k = 1, \dots, K + 1$ with
            \[ Y_{K+1} = \frac{1}{p_K + 1} \,, \quad \text{and} \quad 1 \leq \zeta_k \leq \zeta, \quad \zeta_k \leq \eta_k + 1. \]
    \end{enumerate}}
\end{theorem}

\begin{proof}[Proof of Theorem \ref{thm:order_jump}]
    To show that $p_{K+1} = p_K + 2$ assuming $p_K + 2 \leq p$, we need to prove for all $\tau \in T$ with $|\tau| \leq p_K + 2$ that
    \begin{equation}
        \label{eq:order_jump_main}
        \reva{\alpha^{[K+1]}(\tau) = \beta(\tau) \,, \quad \text{where} \quad \beta(\tau) = 1/\tau! \,.}
    \end{equation}
    This follows for all trees $\tau \in T$ with $|\tau| \leq p_K + 1$ due to Theorem \ref{thm:order_per_sweep}. It remains to show that the same is true for all trees $\tau$ of size $|\tau| = p_K + 2$. 
    \reva{Given $\tau = \reva{\{\tau_1 \cdots \tau_n\}}$ of size $p_K + 2$, we apply Lemma \ref{lemma:ai_gen_form},}
    \[ \alpha^{[K+1]}(\tau) = \frac{1}{|\tau|} \reva{\omega_{K+1}(\tau_1) \cdots \omega_{K+1}(\tau_n)}, \]
    and it remains to show that $\omega_{K+1} (\tau_i) = 1/\tau_i!$ for $i= 1, \dots, n$.
    If $n > 1$, then all trees $\tau_i$ have size less or equal to $p_K$ and, \reva{thus, following Lemma \ref{lemma:ai_gen_form},
    \[ \omega_{K+1}(\tau_i) = (|\tau_i| + 1) \alpha^{[K+1]} (\{\tau_i\}) = \frac{|\tau_i| + 1}{\{\tau_i\}!} = \frac1{\tau_i!} \,, \]
and \eqref{eq:order_jump_main} is proved}. Assume $n = 1$, then let $\tau_1$ be denoted by $\gamma = \reva{\{ \gamma_1\cdots \gamma_m\}}$. We note that $|\gamma| = p_K + 1$. If $p_K < \eta_{K+1}$, then, from the proof of Lemma \ref{lemma:ai_eta_form}, we have
    \begin{align*}
        \omega_{K+1} (\gamma) &= (\frac{1}{|\gamma|} - W_{K,p_K + 1}) \reva{\omega_K (\gamma_1) \cdots \omega_K (\gamma_m)} \\
        &\quad\quad + W_{K,p_K + 1} \reva{\omega_{K+1} (\gamma_1) \cdots \omega_{K+1} (\gamma_m)} = \frac{1}{|\gamma|} \reva{\omega_{K+1} (\gamma_1) \cdots \omega_{K+1} (\gamma_m)},
    \end{align*}
    If $\eta_{K+1} \leq p_K < \eta_{K+1} + \zeta_{K+1}$, then, according to the proof of Lemma \ref{lemma:ai_gen_form} and depending on the structure of $\gamma$, we either apply Lemma \ref{lemma:ai_eta_form} \reva{as in the previous argument, or we obtain,}
    \begin{align*}
        \omega_{K+1} (\gamma) &= (\frac{1}{|\gamma|} - Y_{K+1}) \reva{\omega_K (\gamma_1) \cdots \omega_K (\gamma_m)} + Y_{K+1} \reva{\omega_{K+1} (\gamma_1) \cdots \omega_{K+1} (\gamma_m)} \\
    &= \frac{1}{|\gamma|} \reva{\omega_{K + 1}(\gamma_1) \cdots \omega_{K + 1}(\gamma_m)} .
    \end{align*}
    Since \reva{$|\gamma_i| \leq p_K$ for all $i=1,\dots,m$, we have $\omega_{K+1} (\gamma_i) = 1/\gamma_i!$} and we obtain
    \[ \alpha^{[K+1]} (\tau) = \frac{1}{|\tau|} \frac{1}{|\gamma|} \frac{1}{\gamma_1! \cdots \gamma_m!} = \frac{1}{\tau!}, \]
    and \eqref{eq:order_jump_main} is proved.
\end{proof}

\end{document}
